\documentclass[10pt,a4paper]{amsart}
\usepackage[margin=19mm]{geometry}
\usepackage{amsmath,amssymb,mathtools}
\usepackage[scr=rsfs,cal=euler]{mathalfa}
\usepackage{xfrac}
\usepackage[unicode,hidelinks,bookmarks=false]{hyperref}
\newcommand{\ie}{i.e.\ }
\newcommand{\cf}{cf.\ }
\newcommand{\resp}{respectively\ }
\newcommand{\Subsection}{\subsection}
\providecommand{\nobreakdash}{\nobreak\hbox{-}}
\setlength{\emergencystretch}{2em}
\multlinegap0pt
\usepackage{amssymb}
\usepackage[all]{xy}
\usepackage{bm}
\newtheorem{theorem}{Theorem}
\newtheorem{corollary}{Corollary}
\newtheorem{lemma}{Lemma}
\newcommand{\C}{\mathbb{C}} 
\newcommand{\Fl}{\operatorname{Fl}} 
\newcommand{\FlG}{\operatorname{Fl}} 
\newcommand{\FlT}{\operatorname{Fl}_{\operatorname{ab}}} 
\newcommand{\W}{\mathbb{W}}
\newcommand{\bbS}{\mathbb{S}}
\newcommand{\cL}{\mathcal{L}}
\newcommand{\gt}{{\mathfrak{g}/\mathfrak{t}}}
\newcommand{\loc}{{\operatorname{loc}}} 
\def\parfrac#1#2{\frac{\partial #1}{\partial #2}} 
\title{A mirror theorem for partial flag bundles}
\author{Ionut Ciocan-Fontanine, Yuki Koto}
\date{Source: arXiv:2602.09602v1 (CC BY 4.0)}
\begin{document}
\maketitle

\noindent\emph{Source and licence.} A mirror theorem for partial flag bundles. Version arXiv:2602.09602v1 (CC BY 4.0), distributed by arXiv under the Creative Commons Attribution 4.0 International licence, http://creativecommons.org/licenses/by/4.0/, as recorded in the arXiv licence metadata for this exact version and shown on the abstract page https://arxiv.org/abs/2602.09602v1. Full licence notice: \texttt{licenses/arxiv-2602.09602-CC-BY-4.0.txt}.\par\smallskip

\noindent\textit{Editorial scope.} Counted unit: the article's lemma flag (source0.tex lines 959--961). The article's complete original proof of it is delivered with it (source0.tex lines 964--1003, one \texttt{proof} environment of 40 lines). No mathematics is written, repaired or re-derived: the statement and the proof are the article's own lines, in the article's own notation, and no retained line is substituted or re-flowed.  Retained uncounted as the source-local context the proof needs: lines 751--753; lines 757--760; lines 945--947.  Nothing was omitted from any retained span, so the package carries no omission record.  No retained line contains a citation, so the wrapper carries no bibliography and no bibliography entry.  No retained line contains a figure environment, an image-inclusion command or drawing code, so no image is delivered and the package declares no asset dependency.  Editorial formatting only, in addition: an AMS \texttt{amsart} preamble that restates the article's own declarations and macros used by the retained lines, namely the environments \texttt{theorem}, \texttt{corollary}, \texttt{lemma} on separate counters, together with the standard packages those lines need.  The retained environments print in this document as Theorem 1, Corollary 1, Lemma 1 in order of appearance, whereas the article prints the same items with the numbering of its own full document.  The complete native source of the article as distributed in the arXiv e-print for this exact version is bundled unchanged as \texttt{sources/194329/source0.tex}.
\begin{equation}\label{eqn:Divisor_Equation}
    z\parfrac{I}{t_i^{(m)}}(t_\bullet^\bullet,x,z) = zq_i^{(m)}\parfrac{I}{q_i^{(m)}}(t_\bullet^\bullet,x,z) + H_i^{(m)}\cdot I(t_\bullet^\bullet,x,z).
\end{equation}

\begin{theorem}\label{thm:abelian/nonabelian_correspondence}
    Let $z\cdot I(t^\bullet_\bullet,x,z) \in \cL_{B\times\FlT,\bbS}$ be a point satisfying the Divisor Equation whose specialization $I(t^\bullet,x,z)$ is $\W$-invariant.
    Then, $z \cdot (I(t^\bullet,x,z))_\gt$ lies in $\cL_{B\times\Fl,\bbS}$.
\end{theorem}

\begin{corollary}\label{cor:CFKS}
    For any $f\in H^*_\bbS(\FlT)^{-\W}_\loc[z][\![q^\bullet,\sigma]\!]$ with $f|_{q=\sigma=0}=\omega$, the function $(z\cdot M_{\FlT,\bbS}(\sigma)f) \cup \omega^{-1}$ lies in $\cL_{\FlG,\bbS}$.
\end{corollary}

\begin{lemma} \label{lemma:ab/nonab for flag}
Theorem~\ref{thm:abelian/nonabelian_correspondence} holds in the case when $B$ is a point. 
\end{lemma}

\begin{proof} 
Let $I(q^\bullet_\bullet,t_\bullet^\bullet,x,z)$ be an $H^*_\bbS(\FlT)$-valued formal function with the following properties:
\begin{itemize}
    \item 
    $I(q^\bullet_\bullet,t_\bullet^\bullet,x,z)$ satisfies the Divisor Equation~\eqref{eqn:Divisor_Equation};
    \item 
    $z\cdot I(q^\bullet_\bullet,t_\bullet^\bullet,x,z)$ lies in $\cL_{\FlT,\bbS}$;
    \item 
    the specialization $I(q^\bullet_\bullet,t^\bullet_\bullet,x,z)|_{t_\bullet^\bullet\to t^\bullet}$ is $\W$-invariant.
\end{itemize}
The first and third properties imply that 
\begin{equation}\label{eqn:z_partial_omega_I}
    (I(q^\bullet_\bullet,t^\bullet,x,z))_\gt = (z\partial_\omega I) ((-1)^{r_\bullet-1}q^\bullet,t^\bullet,x,z) \cup \omega^{-1}
\end{equation}
where $z\partial_\omega:=\prod_{\alpha\in\Phi_+} z\partial_{c_1(L_\alpha)}$ and $\partial_{c_1(L_\alpha)}$ denotes the $\C$-linear combination of $\partial_{t_\bullet^\bullet}$ corresponding to $c_1(L_\alpha)$, i.e., the differential operator in $t_\bullet^\bullet$ satisfying 
\[
\partial_{c_1(L_\alpha)} \left( \sum_{m=1}^l \sum_{i=1}^{r_m} t^{(m)}_i H^{(m)}_i \right) = c_1(L_\alpha).
\]
Meanwhile, the second property of $I$ implies that there exist 
\begin{align*}
    \sigma(q^\bullet_\bullet,t^\bullet_\bullet,x) &\in H^*_\bbS(\FlT)_\loc[\![q^\bullet_\bullet,t^\bullet_\bullet,x]\!], \\
    f(q^\bullet_\bullet,t^\bullet_\bullet,x,z) &\in H^*_\bbS(\FlT)_\loc[z][\![q^\bullet_\bullet,t^\bullet_\bullet,x]\!]
\end{align*}
with $\sigma(0,0,0)=0$, $f(0,0,0,z)=1$ such that
\[
I(q^\bullet_\bullet,t_\bullet^\bullet,x,z) = M_{\FlT,\bbS}(\sigma(q^\bullet_\bullet,t^\bullet_\bullet,x)) f(q^\bullet_\bullet,t^\bullet_\bullet,x,z).
\]
The third property of $I$ shows that $\sigma((-1)^{r_\bullet-1}q^\bullet,t^\bullet,x)$ is $\W$-invariant and $f((-1)^{r_\bullet-1}q^\bullet,t^\bullet,x,z)$ is $\W$-invariant.
Since $z\partial_{t^{(m)}_i}\circ M_{\FlT,\bbS}(\sigma) = M_{\FlT,\bbS}(\sigma) \circ (\sigma^*\nabla)_{z\partial_{t^{(m)}_i}}$, the right-hand side of \eqref{eqn:z_partial_omega_I} equals
\[
(M_{\FlT,\bbS}(\sigma((-1)^{r_\bullet-1}q^\bullet,t^\bullet,x)) g((-1)^{r_\bullet-1}q^\bullet,t^\bullet,x,z))\cup\omega^{-1}
\]
for some $g((-1)^{r_\bullet-1}q^\bullet,t^\bullet,x,z)\in H^*_\bbS(\FlT)^{-W}_\loc[z][\![q^\bullet,t^\bullet,x]\!]$ satisfying
$g(0,0,0,z)=\omega$.
Finally, Corollary~\ref{cor:CFKS} yields that 
\[
(I(q^\bullet_\bullet,t^\bullet,x,z))_\gt=(M_{\FlT,\bbS}(\sigma((-1)^{r_\bullet-1}q^\bullet,t^\bullet,x)) g((-1)^{r_\bullet-1}q^\bullet,t^\bullet,x,z))\cup\omega^{-1}
\]
represents a point in $\cL_{\FlG,\bbS}$.
\end{proof}

\end{document}
